% answers.tex —— 2013 年·广东卷·文科 · 参考答案与详细解析片段
% 由 tools/generate_projects.py 依来源总账自 2013/2013_数学真题_广东文_试题.md + 2013/2013_数学真题_广东文_解析正文.md + 2013/2013_数学真题_广东文_答案要点.md 生成，勿手改（改源 MD 后重跑）

\begin{answer}
\gksolhead{1}{}

\textbf{题目}：

设集合 $S=\{x\mid x^2+2x=0,\ x\in\mathbf R\}$，$T=\{x\mid x^2-2x=0,\ x\in\mathbf R\}$，则 $S\cap T=$\mbox{（　　）}

（A）$\{0\}$；（B）$\{0,2\}$；（C）$\{-2,0\}$；（D）$\{-2,0,2\}$。

\Needspace{4\baselineskip}
\textbf{解析}
\Needspace{5\baselineskip}
由 $x^2+2x=0$ 得 $x=0$ 或 $x=-2$，故 $S=\{-2,0\}$；由 $x^2-2x=0$ 得 $x=0$ 或 $x=2$，故 $T=\{0,2\}$。于是
\nopagebreak[4]
\[S\cap T=\{-2,0\}\cap\{0,2\}=\{0\}.\]
选 A。

\gkanswer{答案：A。}

\end{answer}
\begin{answer}
\gksolhead{2}{}

\textbf{题目}：

函数 $y=\dfrac{\lg(x+1)}{x-1}$ 的定义域是\mbox{（　　）}

（A）$(-1,+\infty)$；（B）$[-1,+\infty)$；（C）$(-1,1)\cup(1,+\infty)$；（D）$[-1,1)\cup(1,+\infty)$。

\Needspace{4\baselineskip}
\textbf{解析}
\Needspace{5\baselineskip}
须真数大于零且分母不为零：
\nopagebreak[4]
\[\begin{cases}x+1>0,\\ x-1\neq0,\end{cases}\quad\text{\text{即}}\quad x>-1\ \text{\text{且}}\ x\neq1.\]
故定义域为 $(-1,1)\cup(1,+\infty)$。选 C。

\gkanswer{答案：C。}

\end{answer}
\begin{answer}
\gksolhead{3}{}

\textbf{题目}：

若 $\mathrm i(x+y\mathrm i)=3+4\mathrm i$，$x$，$y\in\mathbf R$，则复数 $x+y\mathrm i$ 的模是\mbox{（　　）}

（A）$2$；（B）$3$；（C）$4$；（D）$5$。

\Needspace{4\baselineskip}
\textbf{解析}
\Needspace{5\baselineskip}
由 $\mathrm i(x+y\mathrm i)=\mathrm i x-y=-y+x\mathrm i=3+4\mathrm i$，比较实部、虚部得 $x=4$，$y=-3$，即 $x+y\mathrm i=4-3\mathrm i$，其模为
\nopagebreak[4]
\[|4-3\mathrm i|=\sqrt{4^2+(-3)^2}=5.\]
选 D。

\gkanswer{答案：D。}

\end{answer}
\begin{answer}
\gksolhead{4}{}

\textbf{题目}：

已知 $\sin\left(\dfrac{5\pi}2+\alpha\right)=\dfrac15$，那么 $\cos\alpha=$\mbox{（　　）}

（A）$-\dfrac25$；（B）$-\dfrac15$；（C）$\dfrac15$；（D）$\dfrac25$。

\Needspace{4\baselineskip}
\textbf{解析}
\Needspace{5\baselineskip}
$\dfrac{5\pi}2=2\pi+\dfrac\pi2$，故
\nopagebreak[4]
\[\sin\left(\frac{5\pi}2+\alpha\right)=\sin\left(\frac\pi2+\alpha\right)=\cos\alpha=\frac15.\]
选 C。

\gkanswer{答案：C。}

\end{answer}
\begin{answer}
\gksolhead{5}{}

\textbf{题目}：

执行如图所示的程序框图，若输入 $n$ 的值为 $3$，则输出 $s$ 的值是\mbox{（　　）}
\nopagebreak[4]
\begin{center}% 2013 广东卷文科第 5 题：程序框图（输入 n，累加 s=s+(i-1)）。
\begin{tikzpicture}[font=\footnotesize, >=Stealth, line width=0.4pt,
  node distance=3.2mm,
  box/.style={draw, minimum height=1.65em, inner xsep=6pt},
  term/.style={draw, rounded corners=6pt, minimum height=1.65em, inner xsep=9pt},
  io/.style={draw, trapezium, trapezium left angle=70, trapezium right angle=110,
             minimum height=1.65em, inner xsep=4pt},
  dec/.style={draw, diamond, aspect=2.5, inner sep=1.2pt}]
  \node[term] (start) {开始};
  \node[io, below=5mm of start] (input) {输入 $n$};
  \node[box, below=5mm of input] (init) {$i=1$，$s=1$};
  \node[dec, below=6mm of init] (cond) {$i\leqslant n$};
  \node[box, below=8mm of cond] (acc) {$s=s+(i-1)$};
  \node[box, below=5mm of acc] (inc) {$i=i+1$};
  \node[io, right=13mm of cond] (output) {输出 $s$};
  \node[term, below=6mm of output] (end) {结束};
  \draw[->] (start) -- (input);
  \draw[->] (input) -- (init);
  \draw[->] (init) -- (cond);
  \draw[->] (cond) -- node[right] {是} (acc);
  \draw[->] (acc) -- (inc);
  \draw[->] (inc.west) -- ++(-1.15,0) |- (cond.west);
  \draw[->] (cond.east) -- node[above] {否} (output.west);
  \draw[->] (output) -- (end);
\end{tikzpicture}
\end{center}
\nopagebreak[4]
（A）$1$；（B）$2$；（C）$4$；（D）$7$。

\Needspace{4\baselineskip}
\textbf{解析}
初值 $i=1$，$s=1$，输入 $n=3$。循环体执行三次：

\begin{itemize}
\item $i=1$：$s=1+(1-1)=1$，$i=2$；
\item $i=2$：$s=1+(2-1)=2$，$i=3$；
\item $i=3$：$s=2+(3-1)=4$，$i=4$。
\end{itemize}

此时 $i=4\leqslant3$ 不成立，跳出循环，输出 $s=4$。选 C。

\gkanswer{答案：C。}

\end{answer}
\begin{answer}
\gksolhead{6}{}

\textbf{题目}：

某三棱锥的三视图如图所示，则该三棱锥的体积是\mbox{（　　）}
\nopagebreak[4]
\begin{center}% 2013 广东卷文科第 6 题：三棱锥的三视图（正视图、侧视图、俯视图）。
\begin{tikzpicture}[font=\scriptsize, >=Stealth, line width=0.4pt,
  x=0.95cm, y=0.95cm]
  % 正视图：直角三角形，竖直直角边在左，高 2、底 1
  \begin{scope}
    \draw (0,0) -- (1,0) -- (0,2) -- cycle;
    \draw[<->] (-0.42,0) -- (-0.42,2) node[midway,left] {$2$};
    \draw[<->] (0,-0.38) -- (1,-0.38) node[midway,below] {$1$};
    \node at (0.5,-1.15) {正视图};
  \end{scope}
  % 侧视图：直角三角形，竖直直角边在右，高 2、底 1
  \begin{scope}[xshift=3.05cm]
    \draw (0,0) -- (1,0) -- (1,2) -- cycle;
    \draw[<->] (0,-0.38) -- (1,-0.38) node[midway,below] {$1$};
    \node at (0.5,-1.15) {侧视图};
  \end{scope}
  % 俯视图：等腰直角三角形，两条直角边均为 1
  \begin{scope}[yshift=-3.55cm]
    \draw (0,0) -- (1,0) -- (1,1) -- cycle;
    \node at (0.5,-1.05) {俯视图};
  \end{scope}
\end{tikzpicture}
\end{center}
\nopagebreak[4]
（A）$\dfrac16$；（B）$\dfrac13$；（C）$\dfrac23$；（D）$1$。

\Needspace{4\baselineskip}
\textbf{解析}
\Needspace{5\baselineskip}
由三视图，正视图、侧视图都是直角边长分别为 $1$ 与 $2$ 的直角三角形，俯视图是两条直角边长均为 $1$ 的等腰直角三角形，故该三棱锥的底面是直角边长均为 $1$ 的等腰直角三角形，高为 $2$。于是
\nopagebreak[4]
\[V=\frac13\cdot\left(\frac12\cdot1\cdot1\right)\cdot2=\frac13.\]
（复核：$\frac13\cdot\frac12\cdot2=\frac13$，与选项 B 相符。）选 B。

\gkanswer{答案：B。}

\end{answer}
\begin{answer}
\gksolhead{7}{}

\textbf{题目}：

垂直于直线 $y=x+1$ 且与圆 $x^2+y^2=1$ 相切于第一象限的直线方程是\mbox{（　　）}

（A）$x+y-\sqrt2=0$；（B）$x+y+1=0$；（C）$x+y-1=0$；（D）$x+y+\sqrt2=0$。

\Needspace{4\baselineskip}
\textbf{解析}
\Needspace{5\baselineskip}
与 $y=x+1$ 垂直，故所求直线斜率为 $-1$，设其为 $x+y=t$。切于第一象限，故切点在 $x$ 轴、$y$ 轴正方向一侧，$t>0$。由圆心到直线的距离等于半径：
\nopagebreak[4]
\[\frac{|t|}{\sqrt2}=1\ \Longrightarrow\ t=\sqrt2\quad(t>0).\]
切点为 $\left(\frac{\sqrt2}2,\frac{\sqrt2}2\right)$，确在第一象限。故所求方程为 $x+y-\sqrt2=0$。选 A。

\gkanswer{答案：A。}

\end{answer}
\begin{answer}
\gksolhead{8}{}

\textbf{题目}：

设 $l$ 为直线，$\alpha$，$\beta$ 是两个不同的平面，下列命题中正确的是\mbox{（　　）}

（A）若 $l\parallel\alpha$，$l\parallel\beta$，则 $\alpha\parallel\beta$；（B）若 $l\perp\alpha$，$l\perp\beta$，则 $\alpha\parallel\beta$；（C）若 $l\perp\alpha$，$l\parallel\beta$，则 $\alpha\parallel\beta$；（D）若 $\alpha\perp\beta$，$l\parallel\alpha$，则 $l\perp\beta$。

\Needspace{4\baselineskip}
\textbf{解析}
（A）反例：$l\parallel\alpha$ 且 $l\parallel\beta$ 时，$\alpha$ 与 $\beta$ 可以相交，故错。
（B）垂直于同一条直线的两个平面平行，正确。
（C）$l\perp\alpha$ 且 $l\parallel\beta$ 时，$l$ 的方向是 $\alpha$ 的法向，又在 $\beta$ 内可取到与 $l$ 平行的直线，故 $\alpha\perp\beta$（而非平行），错。
（D）反例：$\alpha\perp\beta$ 且 $l\parallel\alpha$ 时，$l$ 可能平行于 $\beta$ 内的某条斜线，未必垂直于 $\beta$，错。
选 B。

\gkanswer{答案：B。}

\end{answer}
\begin{answer}
\gksolhead{9}{}

\textbf{题目}：

已知中心在原点的椭圆 $C$ 的右焦点为 $F(1,0)$，离心率等于 $\dfrac12$，则 $C$ 的方程是\mbox{（　　）}

（A）$\dfrac{x^2}3+\dfrac{y^2}4=1$；（B）$\dfrac{x^2}4+\dfrac{y^2}{\sqrt3}=1$；（C）$\dfrac{x^2}4+\dfrac{y^2}2=1$；（D）$\dfrac{x^2}4+\dfrac{y^2}3=1$。

\Needspace{4\baselineskip}
\textbf{解析}
\Needspace{5\baselineskip}
中心在原点、右焦点为 $F(1,0)$ 说明焦点在 $x$ 轴上，且 $c=1$；由 $e=\dfrac ca=\dfrac12$ 得 $a=2$，从而
\nopagebreak[4]
\[b^2=a^2-c^2=4-1=3.\]
故 $C:\dfrac{x^2}4+\dfrac{y^2}3=1$。选 D。

\gkanswer{答案：D。}

\end{answer}
\begin{answer}
\gksolhead{10}{}

\textbf{题目}：

设 $\boldsymbol a$ 是已知的平面向量且 $\boldsymbol a\neq\boldsymbol 0$。关于向量 $\boldsymbol a$ 的分解，有如下四个命题：

① 给定向量 $\boldsymbol b$，总存在向量 $\boldsymbol c$，使 $\boldsymbol a=\boldsymbol b+\boldsymbol c$；

② 给定向量 $\boldsymbol b$ 和 $\boldsymbol c$，总存在实数 $\lambda$ 和 $\mu$，使 $\boldsymbol a=\lambda\boldsymbol b+\mu\boldsymbol c$；

③ 给定单位向量 $\boldsymbol b$ 和正数 $\mu$，总存在单位向量 $\boldsymbol c$ 和实数 $\lambda$，使 $\boldsymbol a=\lambda\boldsymbol b+\mu\boldsymbol c$；

④ 给定正数 $\lambda$ 和 $\mu$，总存在单位向量 $\boldsymbol b$ 和单位向量 $\boldsymbol c$，使 $\boldsymbol a=\lambda\boldsymbol b+\mu\boldsymbol c$。

上述命题中的向量 $\boldsymbol b$，$\boldsymbol c$ 和 $\boldsymbol a$ 在同一平面内且两两不共线，则真命题的个数是\mbox{（　　）}

（A）$1$；（B）$2$；（C）$3$；（D）$4$。

\Needspace{4\baselineskip}
\textbf{解析}
① 取 $\boldsymbol c=\boldsymbol a-\boldsymbol b$ 即可；又 $\boldsymbol a$，$\boldsymbol b$ 不共线时 $\boldsymbol a-\boldsymbol b$ 与两者都不共线，符合题设，真。
② $\boldsymbol b$，$\boldsymbol c$ 不共线，是平面内的一组基底，而 $\boldsymbol a$ 与它们共面，故 $\lambda$，$\mu$ 存在且唯一，真。
③ 由 $\boldsymbol a=\lambda\boldsymbol b+\mu\boldsymbol c$ 及 $|\boldsymbol c|=1$ 得 $|\boldsymbol a-\lambda\boldsymbol b|=\mu$。而 $\boldsymbol a$ 到直线 $\{\lambda\boldsymbol b\}$ 的距离 $d>0$（因 $\boldsymbol a$ 与 $\boldsymbol b$ 不共线），当 $0<\mu<d$ 时方程无解，假。
④ 由 $\boldsymbol a=\lambda\boldsymbol b+\mu\boldsymbol c$ 得 $|\lambda-\mu|\leqslant|\boldsymbol a|\leqslant\lambda+\mu$。当 $|\boldsymbol a|>\lambda+\mu$ 时无解，假。
真命题共 $2$ 个。选 B。

\gkanswer{答案：B。}

\end{answer}
\begin{answer}
\gksolhead{11}{}

\textbf{题目}：

设数列 $\{a_n\}$ 是首项为 $1$，公比为 $-2$ 的等比数列，则 $a_1+|a_2|+a_3+|a_4|=$\gkblank{}。

\Needspace{4\baselineskip}
\textbf{解析}
\Needspace{5\baselineskip}
$a_1=1$，$a_2=1\times(-2)=-2$，$a_3=1\times(-2)^2=4$，$a_4=1\times(-2)^3=-8$。故
\nopagebreak[4]
\[a_1+|a_2|+a_3+|a_4|=1+2+4+8=15.\]

\gkanswer{答案：$15$。}

\end{answer}
\begin{answer}
\gksolhead{12}{}

\textbf{题目}：

若曲线 $y=ax^2-\ln x$ 在点 $(1,a)$ 处的切线平行于 $x$ 轴，则 $a=$\gkblank{}。

\Needspace{4\baselineskip}
\textbf{解析}
\nopagebreak[4]
\[y'=2ax-\frac1x,\]
\Needspace{5\baselineskip}
在 $x=1$ 处切线平行于 $x$ 轴，即 $y'(1)=0$：
\nopagebreak[4]
\[2a-1=0\ \Longrightarrow\ a=\frac12.\]
此时切点为 $(1,a)=\left(1,\frac12\right)$，与题设一致。

\gkanswer{答案：$\dfrac12$。}

\end{answer}
\begin{answer}
\gksolhead{13}{}

\textbf{题目}：

已知变量 $x$，$y$ 满足约束条件 $\begin{cases}x-y+3\geqslant0,\\-1\leqslant x\leqslant1,\\y\geqslant1,\end{cases}$ 则 $z=x+y$ 的最大值是\gkblank{}。

\Needspace{4\baselineskip}
\textbf{解析}
\Needspace{5\baselineskip}
由 $x-y+3\geqslant0$ 得 $y\leqslant x+3$，结合 $-1\leqslant x\leqslant1$、$y\geqslant1$，可行域是以 $(-1,1)$，$(1,1)$，$(1,4)$，$(-1,2)$ 为顶点的四边形。$z=x+y$ 在可行域上取最大值于 $x$ 最大、$y$ 最大之处，即 $(1,4)$：
\nopagebreak[4]
\[z_{\max}=1+4=5.\]

\gkanswer{答案：$5$。}

\end{answer}
\begin{answer}
\gksolhead{14}{}

\textbf{题目}：

已知曲线 $C$ 的极坐标方程为 $\rho=2\cos\theta$，以极点为原点，极轴为 $x$ 轴的正半轴建立直角坐标系，则曲线 $C$ 的参数方程为\gkblank{}。

\Needspace{4\baselineskip}
\textbf{解析}
\Needspace{5\baselineskip}
由 $\rho=2\cos\theta$ 得 $\rho^2=2\rho\cos\theta$，即
\nopagebreak[4]
\[x^2+y^2=2x\ \Longleftrightarrow\ (x-1)^2+y^2=1.\]
\Needspace{5\baselineskip}
这是以 $(1,0)$ 为圆心、$1$ 为半径的圆，故参数方程为
\nopagebreak[4]
\[\begin{cases}x=1+\cos\theta,\\ y=\sin\theta,\end{cases}\quad(\theta\ \text{\text{为参数}}).\]
（也可写成 $x=2\cos^2\theta$，$y=2\sin\theta\cos\theta$。）

\gkanswer{答案：$\begin{cases}x=1+\cos\theta,\\y=\sin\theta\end{cases}$（$\theta$ 为参数，$\theta\in\mathbf R$）。}

\end{answer}
\begin{answer}
\gksolhead{15}{}

\textbf{题目}：

如图，在矩形 $ABCD$ 中，$AB=\sqrt3$，$BC=3$，$BE\perp AC$，垂足为 $E$，则 $ED=$\gkblank{}。
\nopagebreak[4]
\begin{center}% 2013 广东卷文科第 15 题：矩形 ABCD，AB=sqrt3，BC=3，BE 垂直 AC 于 E。
\begin{tikzpicture}[font=\small, >=Stealth, line width=0.4pt]
  \coordinate (A) at (0,0);
  \coordinate (B) at (0,2.2517);
  \coordinate (C) at (3.9,2.2517);
  \coordinate (D) at (3.9,0);
  \coordinate (E) at (0.975,0.5629);
  \draw (A) -- (B) -- (C) -- (D) -- cycle;
  \draw (A) -- (C);
  \draw (B) -- (E);
  \draw (E) -- (D);
  \node[below left] at (A) {$A$};
  \node[above left] at (B) {$B$};
  \node[above right] at (C) {$C$};
  \node[below right] at (D) {$D$};
  \node[above right] at (E) {$E$};
\end{tikzpicture}
\end{center}
\nopagebreak[4]
\Needspace{4\baselineskip}
\textbf{解析}
\Needspace{5\baselineskip}
以 $A$ 为原点、$AD$ 为 $x$ 轴建立坐标系，则
\nopagebreak[4]
\[A(0,0),\quad B(0,\sqrt3),\quad C(3,\sqrt3),\quad D(3,0).\]
\Needspace{5\baselineskip}
$\overrightarrow{AC}=(3,\sqrt3)$，$|AC|=\sqrt{9+3}=2\sqrt3$。$E$ 是 $B$ 在 $AC$ 上的投影，
\nopagebreak[4]
\[AE=\frac{\overrightarrow{AB}\cdot\overrightarrow{AC}}{|AC|}=\frac{0\cdot3+\sqrt3\cdot\sqrt3}{2\sqrt3}=\frac{\sqrt3}2,\]
\Needspace{5\baselineskip}
故 $\overrightarrow{AE}=\dfrac{\sqrt3}{2}\cdot\dfrac{(3,\sqrt3)}{2\sqrt3}=\left(\dfrac34,\dfrac{\sqrt3}4\right)$，即 $E\left(\dfrac34,\dfrac{\sqrt3}4\right)$。于是
\nopagebreak[4]
\[ED=\sqrt{\left(3-\frac34\right)^2+\left(0-\frac{\sqrt3}4\right)^2}=\sqrt{\frac{81}{16}+\frac3{16}}=\sqrt{\frac{84}{16}}=\frac{\sqrt{21}}2.\]
（复核：在 $\triangle AED$ 中 $AE=\frac{\sqrt3}2$，$AD=3$，$\angle EAD=30^\circ$，由余弦定理 $ED^2=\frac34+9-2\cdot\frac{\sqrt3}2\cdot3\cdot\frac{\sqrt3}2=\frac{21}4$，一致。）

\gkanswer{答案：$\dfrac{\sqrt{21}}2$。}

\end{answer}
\begin{answer}
\gksolhead{16}{}

\textbf{题目}：

已知函数 $f(x)=\sqrt2\cos\left(x-\dfrac{\pi}{12}\right)$，$x\in\mathbf R$。

（1） 求 $f\left(\dfrac{\pi}3\right)$ 的值；

（2） 若 $\cos\theta=\dfrac35$，$\theta\in\left(\dfrac{3\pi}2,2\pi\right)$，求 $f\left(\theta-\dfrac{\pi}6\right)$。

\Needspace{4\baselineskip}
\textbf{解析}
\Needspace{5\baselineskip}
\textbf{（1）}
\nopagebreak[4]
\[f\left(\frac\pi3\right)=\sqrt2\cos\left(\frac\pi3-\frac{\pi}{12}\right)=\sqrt2\cos\frac\pi4=\sqrt2\cdot\frac{\sqrt2}2=1.\]

\Needspace{5\baselineskip}
\textbf{（2）} 由 $\theta\in\left(\dfrac{3\pi}2,2\pi\right)$ 得 $\sin\theta<0$，故
\nopagebreak[4]
\[\sin\theta=-\sqrt{1-\cos^2\theta}=-\sqrt{1-\frac9{25}}=-\frac45.\]
\Needspace{5\baselineskip}
于是
\nopagebreak[4]
\[f\left(\theta-\frac\pi6\right)=\sqrt2\cos\left(\theta-\frac\pi6-\frac{\pi}{12}\right)=\sqrt2\cos\left(\theta-\frac\pi4\right),\]
\Needspace{5\baselineskip}
而
\nopagebreak[4]
\[\sqrt2\cos\left(\theta-\frac\pi4\right)=\sqrt2\left(\cos\theta\cos\frac\pi4+\sin\theta\sin\frac\pi4\right)=\cos\theta+\sin\theta,\]
\Needspace{5\baselineskip}
故
\nopagebreak[4]
\[f\left(\theta-\frac\pi6\right)=\frac35-\frac45=-\frac15.\]

\gkanswer{答案（1）：$1$。}

\gkanswer{答案（2）：$-\dfrac15$。}

\end{answer}
\begin{answer}
\gksolhead{17}{}

\textbf{题目}：

从一批苹果中，随机抽取 $50$ 个，其重量（单位：克）的频数分布表如下：

\begin{center}
\begin{tabular}{lllll}
\toprule
分组（重量） & $[80,85)$ & $[85,90)$ & $[90,95)$ & $[95,100)$ \\
\midrule
频数 & $5$ & $10$ & $20$ & $15$ \\
\bottomrule
\end{tabular}
\end{center}

（1） 根据频数分布表计算苹果的重量在 $[90,95)$ 的频率；

（2） 用分层抽样的方法从重量在 $[80,85)$ 和 $[95,100)$ 的苹果中共抽取 $4$ 个，其中重量在 $[80,85)$ 的有几个？

（3） 在（2）中抽出的 $4$ 个苹果中，任取 $2$ 个，求重量在 $[80,85)$ 和 $[95,100)$ 中各有 $1$ 个的概率。

\Needspace{4\baselineskip}
\textbf{解析}
\Needspace{5\baselineskip}
\textbf{（1）} 重量在 $[90,95)$ 的频数为 $20$，样本容量为 $50$，故频率为
\nopagebreak[4]
\[\frac{20}{50}=0.4.\]

\Needspace{5\baselineskip}
\textbf{（2）} $[80,85)$ 有 $5$ 个，$[95,100)$ 有 $15$ 个，共 $20$ 个。按 $1:3$ 分层抽取 $4$ 个，则 $[80,85)$ 中抽取
\nopagebreak[4]
\[4\times\frac{5}{20}=1\ \text{\text{个}}.\]

\Needspace{5\baselineskip}
\textbf{（3）} 设抽出的 $4$ 个中，$[80,85)$ 的那个为 $a$，$[95,100)$ 的三个为 $b_1$，$b_2$，$b_3$。任取 $2$ 个共有
\nopagebreak[4]
\[\mathrm C_4^2=6\ \text{\text{种}},\]
\Needspace{5\baselineskip}
其中两组各有 $1$ 个的取法为 $ab_1$，$ab_2$，$ab_3$ 共 $3$ 种，故所求概率
\nopagebreak[4]
\[P=\frac{3}{6}=\frac12.\]

\gkanswer{答案（1）：$0.4$。}

\gkanswer{答案（2）：$1$ 个。}

\gkanswer{答案（3）：$\dfrac12$。}

\end{answer}
\begin{answer}
\gksolhead{18}{}

\textbf{题目}：

如图①，在边长为 $1$ 的等边三角形 $ABC$ 中，$D$，$E$ 分别是 $AB$，$AC$ 上的点，$AD=AE$，$F$ 是 $BC$ 的中点，$AF$ 与 $DE$ 交于点 $G$。将 $\triangle ABF$ 沿 $AF$ 折起，得到如图②所示的三棱锥 $A-BCF$，其中 $BC=\dfrac{\sqrt2}2$。
\nopagebreak[4]
\begin{center}% 2013 广东卷文科第 18 题图①：等边三角形 ABC 中的折痕（AF 与 DE 交于 G）。
\begin{tikzpicture}[font=\small, >=Stealth, line width=0.4pt]
  \coordinate (B) at (0,0);
  \coordinate (C) at (3,0);
  \coordinate (A) at (1.5,2.5981);
  \coordinate (F) at (1.5,0);
  \coordinate (D) at (0.5,0.8660);
  \coordinate (E) at (2.5,0.8660);
  \coordinate (G) at (1.5,0.8660);
  \draw (A) -- (B) -- (C) -- cycle;
  \draw (A) -- (F);
  \draw (D) -- (E);
  \fill (G) circle (1.1pt);
  \node[above] at (A) {$A$};
  \node[below left] at (B) {$B$};
  \node[below right] at (C) {$C$};
  \node[below] at (F) {$F$};
  \node[left] at (D) {$D$};
  \node[right] at (E) {$E$};
  \node[above left] at (G) {$G$};
  \node (lbl) at (1.5,-0.85) {图};
  \node[circle, draw, inner sep=0.4pt, font=\scriptsize, right=1pt of lbl] {1};
\end{tikzpicture}
\end{center}
\nopagebreak[4]
\nopagebreak[4]
\begin{center}% 2013 广东卷文科第 18 题图②：折叠后的三棱锥 A-BCF（D、E、G 为折痕上的点）。
\begin{tikzpicture}[font=\small, >=Stealth, line width=0.4pt]
  \coordinate (A) at (1.55,3.60);
  \coordinate (B) at (0,0);
  \coordinate (C) at (3.40,0.55);
  \coordinate (F) at (1.45,0.75);
  \coordinate (D) at (0.5167,1.20);
  \coordinate (E) at (2.7833,1.5667);
  \coordinate (G) at (1.4833,1.70);
  % 被遮挡的棱与辅助线
  \draw[dashed] (A) -- (F);
  \draw[dashed] (F) -- (B);
  \draw[dashed] (F) -- (C);
  \draw[dashed] (G) -- (D);
  \draw[dashed] (G) -- (E);
  % 可见的棱
  \draw (A) -- (B) -- (C) -- (A);
  \draw (D) -- (E);
  \fill (F) circle (1.1pt);
  \fill (G) circle (1.1pt);
  \node[above] at (A) {$A$};
  \node[below left] at (B) {$B$};
  \node[right] at (C) {$C$};
  \node[below] at (F) {$F$};
  \node[left] at (D) {$D$};
  \node[right] at (E) {$E$};
  \node[above left] at (G) {$G$};
  \node (lbl) at (1.55,-0.75) {图};
  \node[circle, draw, inner sep=0.4pt, font=\scriptsize, right=1pt of lbl] {2};
\end{tikzpicture}
\end{center}
\nopagebreak[4]
（1） 证明：$DE\parallel$ 平面 $BCF$；

（2） 证明：$CF\perp$ 平面 $ABF$；

（3） 当 $AD=\dfrac23$ 时，求三棱锥 $F-DEG$ 的体积 $V_{F-DEG}$。

\Needspace{4\baselineskip}
\textbf{解析}
\Needspace{5\baselineskip}
\textbf{（1）} 折叠前 $AD=AE$，$AB=AC=1$，故 $\dfrac{AD}{AB}=\dfrac{AE}{AC}$，从而 $DE\parallel BC$。折叠是绕 $AF$ 的刚体运动，$DE$ 与 $BC$ 的平行关系不变，故折叠后仍有 $DE\parallel BC$。又 $BC\subset$ 平面 $BCF$，$DE\not\subset$ 平面 $BCF$，所以
\nopagebreak[4]
\[DE\parallel\text{\text{平面} }BCF.\]

\Needspace{5\baselineskip}
\textbf{（2）} 折叠前 $AF\perp BC$，而 $F$ 在折痕 $AF$ 上，故折叠后 $AF\perp BF$，$AF\perp CF$ 均保持。又 $BF=CF=\dfrac12$，$BC=\dfrac{\sqrt2}2$，于是
\nopagebreak[4]
\[BF^2+CF^2=\frac14+\frac14=\frac12=BC^2,\]
\Needspace{5\baselineskip}
故 $\angle BFC=90^\circ$，即 $CF\perp BF$。由 $CF\perp AF$，$CF\perp BF$，$AF\cap BF=F$，得
\nopagebreak[4]
\[CF\perp\text{\text{平面} }ABF.\]

\Needspace{5\baselineskip}
\textbf{（3）} 由（2）$CF\perp$ 平面 $ABF$，而 $AF=\dfrac{\sqrt3}2$（边长 $1$ 的等边三角形的中线），故
\nopagebreak[4]
\[V_{A-BCF}=\frac13\cdot S_{\triangle ABF}\cdot CF=\frac13\cdot\left(\frac12\cdot AF\cdot BF\right)\cdot CF=\frac13\cdot\frac12\cdot\frac{\sqrt3}2\cdot\frac12\cdot\frac12=\frac{\sqrt3}{48}.\]
\Needspace{5\baselineskip}
取以 $A$ 为心、比为 $\dfrac23$ 的位似变换：由 $AD=AE=\dfrac23$，$AG=\dfrac23AF$ 知它把 $B$，$C$，$F$ 分别映为 $D$，$E$，$G$，故把平面 $BCF$ 映为平面 $DEG$（两面平行），把三棱锥 $A-BCF$ 映为 $A-DEG$，于是
\nopagebreak[4]
\[V_{A-DEG}=\left(\frac23\right)^3V_{A-BCF}=\frac8{27}\cdot\frac{\sqrt3}{48}.\]
\Needspace{5\baselineskip}
又 $A$ 到平面 $DEG$ 的距离为 $\dfrac23AF$，$F$ 到平面 $DEG$ 的距离为 $\dfrac13AF$（因 $G$ 在 $AF$ 上且 $AG=\dfrac23AF$），两面同底 $DEG$，故
\nopagebreak[4]
\[V_{F-DEG}=\frac12V_{A-DEG}=\frac4{27}\cdot\frac{\sqrt3}{48}=\frac{\sqrt3}{324}.\]
（数值复核：$|BC|=\frac{\sqrt2}2$，$\frac{\sqrt3}{324}\approx0.00535$；亦可取 $F$ 为原点、$FB$，$FC$，$FA$ 为三条两两垂直的轴直接算行列式，结果相同。）

\gkanswer{答案（1）：证明见解析。}

\gkanswer{答案（2）：证明见解析。}

\gkanswer{答案（3）：$\dfrac{\sqrt3}{324}$。}

\end{answer}
\begin{answer}
\gksolhead{19}{}

\textbf{题目}：

设各项均为正数的数列 $\{a_n\}$ 的前 $n$ 项和为 $S_n$，满足 $4S_n=a_{n+1}^2-4n-1$，$n\in\mathbf N^*$，且 $a_2$，$a_5$，$a_{14}$ 构成等比数列。

（1） 证明：$a_2=\sqrt{4a_1+5}$；

（2） 求数列 $\{a_n\}$ 的通项公式；

（3） 证明：对一切正整数 $n$，有 $\dfrac1{a_1a_2}+\dfrac1{a_2a_3}+\cdots+\dfrac1{a_na_{n+1}}<\dfrac12$。

\Needspace{4\baselineskip}
\textbf{解析}
\Needspace{5\baselineskip}
\textbf{（1）} 在 $4S_n=a_{n+1}^2-4n-1$ 中取 $n=1$，得 $4a_1=a_2^2-5$，即 $a_2^2=4a_1+5$。因各项均为正数，$a_2>0$，故
\nopagebreak[4]
\[a_2=\sqrt{4a_1+5}.\]

\Needspace{5\baselineskip}
\textbf{（2）} 当 $n\geqslant2$ 时，$4S_{n-1}=a_n^2-4(n-1)-1=a_n^2-4n+3$。与原式相减得
\nopagebreak[4]
\[4a_n=a_{n+1}^2-a_n^2-4,\]
\Needspace{5\baselineskip}
即 $a_{n+1}^2=(a_n+2)^2$。由 $a_{n+1}>0$，$a_n+2>0$ 得
\nopagebreak[4]
\[a_{n+1}=a_n+2\quad(n\geqslant2),\]
故 $\{a_n\}$ 从第 $2$ 项起是公差为 $2$ 的等差数列，$a_n=a_2+2(n-2)$（$n\geqslant2$）。

\Needspace{5\baselineskip}
由 $a_2$，$a_5$，$a_{14}$ 成等比数列，$a_5=a_2+6$，$a_{14}=a_2+24$，得
\nopagebreak[4]
\[(a_2+6)^2=a_2(a_2+24)\ \Longrightarrow\ 12a_2=36\ \Longrightarrow\ a_2=3.\]
\Needspace{5\baselineskip}
再由（1）$a_1=\dfrac{a_2^2-5}4=\dfrac{9-5}4=1$。于是对 $n\geqslant2$ 有 $a_n=3+2(n-2)=2n-1$，且 $n=1$ 时 $2\times1-1=1=a_1$ 也成立，故
\nopagebreak[4]
\[a_n=2n-1\quad(n\in\mathbf N^*).\]

\Needspace{5\baselineskip}
\textbf{（3）} 由（2），
\nopagebreak[4]
\[\frac1{a_na_{n+1}}=\frac1{(2n-1)(2n+1)}=\frac12\left(\frac1{2n-1}-\frac1{2n+1}\right),\]
\Needspace{5\baselineskip}
故
\nopagebreak[4]
\[\sum_{k=1}^n\frac1{a_ka_{k+1}}=\frac12\left[\left(1-\frac13\right)+\cdots+\left(\frac1{2n-1}-\frac1{2n+1}\right)\right],\]
\Needspace{5\baselineskip}
即
\nopagebreak[4]
\[\sum_{k=1}^n\frac1{a_ka_{k+1}}=\frac12\left(1-\frac1{2n+1}\right)<\frac12.\]

\gkanswer{答案（1）：证明见解析。}

\gkanswer{答案（2）：$a_n=2n-1$。}

\gkanswer{答案（3）：证明见解析。}

\end{answer}
\begin{answer}
\gksolhead{20}{}

\textbf{题目}：

已知抛物线 $C$ 的顶点为原点，其焦点 $F(0,c)$（$c>0$）到直线 $l:x-y-2=0$ 的距离为 $\dfrac{3\sqrt2}2$。设 $P$ 为直线 $l$ 上的点，过点 $P$ 作抛物线 $C$ 的两条切线 $PA$，$PB$，其中 $A$，$B$ 为切点。

（1） 求抛物线 $C$ 的方程；

（2） 当点 $P(x_0,y_0)$ 为直线 $l$ 上的定点时，求直线 $AB$ 的方程；

（3） 当点 $P$ 在直线 $l$ 上移动时，求 $|AF|\cdot|BF|$ 的最小值。

\Needspace{4\baselineskip}
\textbf{解析}
\Needspace{5\baselineskip}
\textbf{（1）} 焦点 $F(0,c)$ 到直线 $x-y-2=0$ 的距离为
\nopagebreak[4]
\[\frac{|0-c-2|}{\sqrt{1^2+(-1)^2}}=\frac{|c+2|}{\sqrt2}=\frac{3\sqrt2}2,\]
\Needspace{5\baselineskip}
故 $|c+2|=3$，即 $c=1$ 或 $c=-5$。由 $c>0$ 得 $c=1$，抛物线为
\nopagebreak[4]
\[x^2=4y.\]

\Needspace{5\baselineskip}
\textbf{（2）} 设 $A(x_1,y_1)$，$B(x_2,y_2)$ 在 $x^2=4y$ 上。由 $y=\dfrac{x^2}4$，抛物线在 $A$ 处的切线为
\nopagebreak[4]
\[y=\frac{x_1}2x-\frac{x_1^2}4,\quad\text{\text{即}}\quad x_1x=2(y+y_1).\]
\Needspace{5\baselineskip}
它过 $P(x_0,y_0)$，故 $x_1x_0=2(y_0+y_1)$，即点 $A$ 满足 $x_0x_1-2y_1-2y_0=0$。同理 $B$ 也满足 $x_0x_2-2y_2-2y_0=0$。所以 $A$，$B$ 都在直线
\nopagebreak[4]
\[x_0x-2y-2y_0=0\]
上，该直线即弦 $AB$ 所在直线（$A\neq B$ 时唯一确定）。

\Needspace{5\baselineskip}
\textbf{（3）} 将 $y=\dfrac{x_0x}2-y_0$ 代入 $x^2=4y$ 得
\nopagebreak[4]
\[x^2-2x_0x+4y_0=0,\]
\Needspace{5\baselineskip}
故 $x_1+x_2=2x_0$，$x_1x_2=4y_0$。由抛物线定义（准线 $y=-1$），$|AF|=y_1+1$，$|BF|=y_2+1$，于是
\nopagebreak[4]
\[|AF|\cdot|BF|=(y_1+1)(y_2+1)=\frac{x_1^2x_2^2}{16}+\frac{x_1^2+x_2^2}4+1.\]
\Needspace{5\baselineskip}
代入 $x_1x_2=4y_0$，$x_1^2+x_2^2=(x_1+x_2)^2-2x_1x_2=4x_0^2-8y_0$，得
\nopagebreak[4]
\[|AF|\cdot|BF|=y_0^2+x_0^2-2y_0+1=x_0^2+(y_0-1)^2.\]
\Needspace{5\baselineskip}
又 $P$ 在 $l$ 上，$x_0=y_0+2$，故
\nopagebreak[4]
\[|AF|\cdot|BF|=(y_0+2)^2+(y_0-1)^2=2y_0^2+2y_0+5=2\left(y_0+\frac12\right)^2+\frac92\geqslant\frac92,\]
当 $y_0=-\dfrac12$，即 $P\left(\dfrac32,-\dfrac12\right)$ 时取等号。故最小值为 $\dfrac92$。

\gkanswer{答案（1）：$x^2=4y$。}

\gkanswer{答案（2）：$x_0x-2y-2y_0=0$。}

\gkanswer{答案（3）：$\dfrac92$。}

\end{answer}
\begin{answer}
\gksolhead{21}{}

\textbf{题目}：

设函数 $f(x)=x^3-kx^2+x$（$x\in\mathbf R$）。

（1） 当 $k=1$ 时，求函数 $f(x)$ 的单调区间；

（2） 当 $k<0$ 时，求函数 $f(x)$ 在 $[k,-k]$ 上的最小值 $m$ 和最大值 $M$。

\Needspace{4\baselineskip}
\textbf{解析}
\Needspace{5\baselineskip}
\textbf{（1）} 当 $k=1$ 时 $f(x)=x^3-x^2+x$，
\nopagebreak[4]
\[f'(x)=3x^2-2x+1=3\left(x-\frac13\right)^2+\frac23>0,\]
故 $f(x)$ 在 $(-\infty,+\infty)$ 上单调递增，单调递增区间为 $(-\infty,+\infty)$，无单调递减区间。

\Needspace{5\baselineskip}
\textbf{（2）} 注意到
\nopagebreak[4]
\[f(x)-k=x^3-kx^2+x-k=x^2(x-k)+(x-k)=(x-k)(x^2+1).\]
\Needspace{5\baselineskip}
当 $k<0$ 时 $[k,-k]$ 上 $x\geqslant k$，故 $f(x)-k\geqslant0$，且仅在 $x=k$ 处取等号，所以
\nopagebreak[4]
\[m=f(k)=k.\]
\Needspace{5\baselineskip}
又
\nopagebreak[4]
\[f(x)-(-2k^3-k)=x^3-kx^2+x+2k^3+k=(x+k)\left(x^2-2kx+2k^2+1\right),\]
\Needspace{5\baselineskip}
而
\nopagebreak[4]
\[x^2-2kx+2k^2+1=(x-k)^2+k^2+1>0,\]
\Needspace{5\baselineskip}
在 $[k,-k]$ 上 $x+k\leqslant0$，故 $f(x)\leqslant-2k^3-k$，且仅在 $x=-k$ 处取等号，所以
\nopagebreak[4]
\[M=f(-k)=-k^3-k\cdot k^2-k=-2k^3-k.\]
（复核：$k=-1$ 时 $f(x)=x^3+x^2+x$ 在 $[-1,1]$ 上递增，$m=f(-1)=-1=k$，$M=f(1)=3=-2k^3-k$；$k=-2$ 时 $m=-2$，$M=f(2)=18=-2(-8)+2$，均相符。）

\gkanswer{答案（1）：单调递增区间为 $(-\infty,+\infty)$，无单调递减区间。}

\gkanswer{答案（2）：$m=k$，$M=-2k^3-k$。}

\end{answer}